\section{Estudio del cambio de base en los funtores homológicos covariantes de módulos}
\label{section:III.7}

\subsection{Funtores de $A$-módulos}
\label{subsection:III.7.1}

\begin{env}[7.1.1]
\label{III.7.1.1}
Dado un anillo $A$, no necesariamente conmutativo, denotemos por
$\cat{Ab}_A$ la categoría de los $A$-módulos a izquierda, y
simplemente por $\cat{Ab}$ la categoría de los $\bb Z$-módulos,
que es la categoría de los grupos abelianos. Sea
\[
  T:\cat{Ab}_A\longrightarrow\cat{Ab}
\]
un funtor covariante aditivo, y sea $M$ un $(A,A)$-bimódulo.
Entonces $T(M)$ está dotado de manera natural de una estructura de
$A$-módulo a derecha.

En efecto, para cada $a\in A$, denotemos por $h_{a,M}$, o
simplemente por $h_a$, el endomorfismo
\[
  x\longmapsto xa
\]
del $A$-módulo a izquierda $M$. Por hipótesis, $T(h_a)$ es un
endomorfismo del $\bb Z$-módulo $T(M)$. Como $T$ es covariante y
aditivo, para $a,b\in A$ se tiene
\[
  T(h_{ab})
  =
  T(h_b\circ h_a)
  =
  T(h_b)\circ T(h_a)
\]
y
\[
  T(h_{a+b})
  =
  T(h_a+h_b)
  =
  T(h_a)+T(h_b).
\]
Por tanto,
\[
  (y,a)\longmapsto T(h_a)(y)
\]
define una multiplicación escalar de $A$-módulo a derecha sobre
$T(M)$. En particular, $T(A_s)$ es un $A$-módulo a derecha.
\end{env}

\begin{env}[7.1.2]
\label{III.7.1.2}
Supongamos que $A$ es conmutativo. De \sref{III.7.1.1} se deduce
que, para todo $A$-módulo $M$, el grupo $T(M)$ es naturalmente un
$A$-módulo. Además, si $u:M\to N$ es un homomorfismo
$A$-lineal, entonces, para todo $a\in A$,
\[
  u\circ h_{a,M}=h_{a,N}\circ u,
\]
y, por consiguiente,
\[
  T(u)\circ T(h_{a,M})
  =
  T(h_{a,N})\circ T(u).
\]
Así, $T(u):T(M)\to T(N)$ es $A$-lineal, y $T$ puede considerarse
como un funtor covariante aditivo de $\cat{Ab}_A$ en sí misma.

Más precisamente, esto define una equivalencia canónica entre la
categoría de los funtores covariantes aditivos
\[
  \cat{Ab}_A\longrightarrow\cat{Ab}
\]
y la categoría de los funtores covariantes $A$-lineales
\[
  T:\cat{Ab}_A\longrightarrow\cat{Ab}_A,
\]
es decir, de los funtores que satisfacen
\[
  T(h_{a,M})=h_{a,T(M)}
  \qquad(a\in A).
\]
El funtor de olvido
\[
  I:\cat{Ab}_A\longrightarrow\cat{Ab},
\]
que asocia a un $A$-módulo su $\bb Z$-módulo subyacente, es exacto
y fiel. Por tanto, los dos funtores correspondientes bajo la
equivalencia precedente tienen las mismas propiedades de exactitud.
\end{env}

\begin{env}[7.1.3]
\label{III.7.1.3}
Sigamos suponiendo que $A$ es conmutativo. Sea $B$ una
$A$-álgebra, no necesariamente conmutativa, y sea
\[
  \rho:A\longrightarrow B
\]
el homomorfismo de anillos que define la estructura de álgebra.
Este define un
\oldpage[III]{176}
funtor covariante aditivo
\[
  \rho_*:M\longmapsto M_{[\rho]}
\]
de la categoría $\cat{Ab}_B$ de los $B$-módulos a izquierda a
$\cat{Ab}_A$. Por composición se obtiene el funtor covariante
aditivo
\[
  T_{(B)}:
  \cat{Ab}_B
  \xrightarrow{\rho_*}
  \cat{Ab}_A
  \xrightarrow{T}
  \cat{Ab}.
\]
Por razones tipográficas, este funtor se denota también por
$T^{(B)}$ o por $T\otimes_AB$, y se dice que se obtiene de $T$ por
extensión de escalares de $A$ a $B$.

Si $B$ es conmutativa, $T_{(B)}$ puede considerarse, por supuesto,
como un funtor de $\cat{Ab}_B$ en sí misma
\sref{III.7.1.2}. Si $B$ es conmutativa y $C$ es una $B$-álgebra,
entonces
\[
  T_{(C)}=(T_{(B)})_{(C)}.
\]
La extensión de escalares es manifiestamente funtorial y aditiva
en $T$. Además, si $T$ conmuta con los límites inductivos o con las
sumas directas, lo mismo ocurre con $T_{(B)}$; y si $T$ es exacto
por la izquierda, exacto por la derecha o exacto,
respectivamente, lo mismo vale para $T_{(B)}$. En efecto,
$\rho_*$ es exacto y conmuta con los límites inductivos y las sumas
directas.
\end{env}

\begin{env}[7.1.4]
\label{III.7.1.4}
Supongamos de nuevo que $A$ es conmutativo, y sea
\[
  T:\cat{Ab}_A\longrightarrow\cat{Ab}_A
\]
un funtor covariante aditivo $A$-lineal que conmuta con los
límites inductivos. Para toda parte multiplicativa $S$ de $A$ y
todo $A$-módulo $M$, existe un isomorfismo canónico funtorial de
$A$-módulos
\begin{env}[7.1.4.1]
\label{III.7.1.4.1}
\[
  T(S^{-1}M)
  \xrightarrow{\sim}
  S^{-1}T(M).
  \tag{7.1.4.1}\label{III.7.1.4.1.eq}
\]
\end{env}

Supongamos primero que $S$ es el conjunto de las potencias $f^n$,
$n\geq0$, de un elemento $f\in A$. Entonces
\[
  M_f=\varinjlim M_n,
\]
donde $(M_n,\vphi_{nm})$ es el sistema inductivo de $A$-módulos
con $M_n=M$ y
\[
  \vphi_{nm}:z\longmapsto f^{\,n-m}z
  \qquad\text{\sref[0]{0.1.6.1}}.
\]
La fórmula \eqref{III.7.1.4.1.eq} se deduce en este caso de la
hipótesis sobre $T$. Para $S$ arbitraria, el módulo $S^{-1}M$ es
el límite inductivo de los $M_f$, para $f\in S$
\sref[0]{0.1.4.5}, y se aplica el mismo razonamiento.

La funtorialidad de \eqref{III.7.1.4.1.eq} muestra además que se
trata de un isomorfismo de $S^{-1}A$-módulos. Así, salvo
isomorfismo canónico, se puede escribir
\begin{env}[7.1.4.2]
\label{III.7.1.4.2}
\[
  T_{(S^{-1}A)}(S^{-1}M)
  =
  S^{-1}T(M)
  =
  T(S^{-1}M).
  \tag{7.1.4.2}\label{III.7.1.4.2.eq}
\]
\end{env}
Si $S=A-\mathfrak p$ es el complementario de un ideal primo
$\mathfrak p$ de $A$, se escribe $T_{\mathfrak p}$ en lugar de
$T_{(A_{\mathfrak p})}$.
\end{env}

\begin{proposition}[7.1.5]
\label{III.7.1.5}
Bajo las hipótesis de \sref{III.7.1.4}, supongamos que
$T_{\mathfrak m}$ es exacto por la izquierda, respectivamente
exacto por la derecha o exacto, para todo ideal maximal
$\mathfrak m$ de $A$. Entonces $T$ es exacto por la izquierda,
respectivamente exacto por la derecha o exacto.
\end{proposition}

\begin{proof}
Si dos submódulos $N$ y $P$ de un $A$-módulo $M$ satisfacen
\[
  N_{\mathfrak m}=P_{\mathfrak m}
\]
para todo ideal maximal $\mathfrak m$ de $A$, entonces $N=P$
(Bourbaki, \emph{Algèbre commutative}, cap.~II, \S~3,
n.º~3, teo.~1).
\end{proof}
